\documentclass[ 		% Clase del documento
	12pt,				% Tamaño de letra
	a4paper,			% Tamaño de papel
	twoside,			% Impresión a doble cara
	openright,			% La apertura de cap. a la dcha.
%	draft       		% Versión borrador (sin figuras)
	final       		% Versión final
]{book}

%--- Márgenes del documento
\usepackage[top=2.5cm,    % Superior
            bottom=2.5cm, % Inferior
            inner=3cm,  % Interior 
            outer=3cm     % Exterior
]{geometry}

\newif\ifspanish\spanishtrue
% EDITAR: Descomentar si se desea que el idioma pral. sea inglés
%\spanishfalse 
\newif\ifpageonfooter\pageonfooterfalse
% EDITAR: Descomentar si se desea número de página en el pie de página
%\pageonfootertrue

\usepackage{uclmTFMhum}
%%
%% Añade aquí los paquetes adicionales que necesites
%%
\usepackage[natbibapa]{apacite}

% -------------------------
% -------------------------
% -------------------------
% BEGIN_FOLD
% -------------------------
% -------------------------
% -------------------------
% EDITAR: Datos del documento. 
% Definición de variables empleadas en el documento por lo que no son
% traducidos. Cuando algún campo puede tener varias líneas aparecen dos
% campos señalados como <campo>Primera y <campo>Segunda. Si no se desea 
% emplear los campos opcionales (OPT.) estos deben comentarse.
% -------------------------
\autor{Autor/a}
\email{autor-a@uclm.es}					
\instEdu{UNIVERSIDAD DE CASTILLA-LA MANCHA}
\centroEdu{FACULTAD DE LETRAS}
\titulacion{MÁSTER UNIVERSITARIO EN INVESTIGACIÓN EN LETRAS Y HUMANIDADES} 
\tipoDoc{TRABAJO FIN DE MÁSTER}
% Si las fechas se desean en inglés hay que ponerla explícita.
\mesTF{julio}        	% Mes de defensa (incluir siempre)
\monthTF{July}        	% En inglés (incluir siempre)
\yearTF{2023}        	% Año de defensa
\cityTF{Ciudad Real}	% Ciudad de defensa



% EDITAR: Para cambiar ajustes y colores de los enlaces y marcadores en el texto

\setcounter{tocdepth}{1} % Suprime del índice los apdos. de nivel mayor
\setcounter{secnumdepth}{3} % Suprime numeración de secciones de nivel mayor

\makeatletter
% Más información en: https://en.wikibooks.org/wiki/LaTeX/Hyperlinks#Customization
\hypersetup{% (Metadatos del fichero PDF generado)
    linktocpage=true,        % true = enlace al nº de pág., false=texto completo
	colorlinks=true,         % true=colorea texto del enlace, false=recuadra el texto
%    hidelinks,               % Oculta colores en los enlaces (en negro)
    linkcolor=UCLMred,       % Color links internos
    anchorcolor=UCLMred,     % Color para anclas a texto
	urlcolor=aquaESI,        % Color para URL enlazadas
	citecolor=UCLMred,       % Color para citas a bibliografía
	bookmarksopen=true,      % Abre PDF con panel de marcadores abierto
	bookmarksnumbered=true,  % Incluye números en marcadores
	pdftoolbar=true,         % Muestra la toolbar de Acrobat
	pdfmenubar=true,	     % Muestra la menubar de Acrobat
	pdftitle={\@title},      % Título
	pdfauthor={\@author},    % Autor
	pdfsubject={\@tipoDoc}   % Tipo de de documento
}
\makeatother















% -------------------------
% -------------------------
% -------------------------
% -------------------------
%
% CUERPO del documento
%
% -------------------------
\begin{document}
%--- PORTADAS + FRONTMATTER
\frontmatter
\pagestyle{empty}  % Páginas sin cabecera ni pies
% OJO: Editar este fichero a voluntad para ajustarlo al resultado deseado.
% OJO: Poner especial atención en usar los adjetivos de género correctos.
% -------------------------
% -------------------------
% -------------------------
% PORTADAS: (Incluidas páginas para créditos y dedicatoria)
% -------------------------
% -------------------------------------------------------------------------
% -------------------------------------------------------------------------
% -------------------------------------------------------------------------
%--- PORTADA TFM LETRAS
\begin{titlepage}
    \begin{center}
        {\includegraphics[width=3cm]{uclm3arcos}}\\[0.4cm]
        {\LARGE UNIVERSIDAD DE CASTILLA-LA MANCHA}
        \vspace{1cm}
        {\par\Large FACULTAD DE LETRAS} 
        \vspace{0.6cm}
    {\par\textsc{MÁSTER UNIVERSITARIO EN INVESTIGACIÓN\\ 
    EN LETRAS Y HUMANIDADES}}
	\vfill
    {\huge\textsc{Título del trabajo (español)}}\par
    \bigskip
    \parbox{.6\linewidth}{\dingline{167}}
    \bigskip
    \par{\Large\textsc{Título del trabajo (inglés)}}
    \par\vspace{2cm}
    \textsc{(Autor/a)}
    \medbreak
    Director/a: Prof. Nombre Apellidos (Director/a)\par
    Depto. del Director/a\par
    \vspace{4cm}
    \noindent%
    Ciudad real, fecha
    \end{center}
    \vfill
    \par\noindent\makebox[2.5in]{\hrulefill} \hfill\makebox[2.5in]{\hrulefill}%
    \par\noindent\makebox[2.5in]{Firma Autor/a}      
    \hfill\makebox[2.5in]{Firma Director/a}%
\end{titlepage}
%---

% OPT.: DEDICATORIA (1 pág. máximo) comentar si no se desea incluir.
% Aunque opcional, no se debería perder la oportunidad de poder 
% dedicar el trabajo a alguien MUY especial. Debe ocupar como mucho dos líneas
% (no confundir con los agradecimientos).
% EDITAR: Dedicatoria.

\null\vspace{\stretch{1}}
\begin{flushright}
\emph{A alguien muy especial\\
Por el motivo que lo hace especial}
\end{flushright}
\vspace{\stretch{2}}\null
\cleardoublepage % Diseño de portadas
%--- Ajustes del documento.
\pagestyle{plain}	% Páginas sólo con numeración inferior al pie

% -------------------------
%
% RESUMEN:
% 
%

% EDITAR: Resumen (máx. 1 pág.)
%\cleardoublepage % Se incluye para modificar el contador de página antes de añadir bookmark
\phantomsection  % Necesario con hyperref
\addcontentsline{toc}{chapter}{Resumen} % Añade al TOC.
\selectlanguage{spanish} % Selección de idioma del resumen.
\makeatletter
\begin{center} %
   {\textsc{TRABAJO FIN DE MÁSTER - FACULTAD DE LETRAS (UCLM)}\par} % Tipo de trabajo
   \vspace{1cm} %  
   {\textbf{\Large Título en español}\par}  % Título del trabajo
   \vspace{0.4cm} %
   {\@autor \\ \@cityTF,{} \@mesTF{} \@yearTF\par} 
   \vspace{0.9cm} %
   {\textbf{\large\textsf{Resumen}}\par} % Título de resumen
\end{center}   
\makeatother %
En una página como máximo, el resumen explicará de modo breve la problemática que trata de resolver el trabajo \emph{(<<Qué>>)}, la metodología para  abordar su solución  \emph{<<Cómo>>)} y los resultados obtenidos. En los trabajos cuyo idioma principal sea el inglés, el orden de \textsf{Resumen} y \textsf{Abstract} se invertirá.



%---
\cleardoublepage % Se incluye para modificar el contador de página antes de añadir 





% EDITAR: Abstract (máx. 1 pág.)
%---
\phantomsection  % Necesario con hyperref
\addcontentsline{toc}{chapter}{Abstract} % Añade al TOC.
\selectlanguage{english} % Selección de idioma del resumen.
\makeatletter
\begin{center} %
   {\textsc{MASTER THESIS - FACULTAD DE LETRAS (UCLM)}\par}
   \vspace{1cm} %  
   {\textbf{\Large Title in English}\par}
   \vspace{0.4cm} %
   {\@autor \\ \@cityTF,{} \@monthTF{} \@yearTF\par} 
   \vspace{0.9cm} %
   {\textbf{\large\textsf{Abstract}}\par} 
\end{center}   
\makeatother %
\emph{English version for the abstract.}
\cleardoublepage % Se incluye para modificar el contador de página antes de añadir 

 % Resumen/Abstract
% Selección del idioma principal del texto
\ifspanish
	\selectlanguage{spanish}
\else
	\selectlanguage{english}
\fi

% EDITAR: Según el gusto del usuario

% -------------------------
%
% AGRADECIMIENTOS (recomendable máx. 1 pág.)
%
% -------------------------
\phantomsection % Necesario con hyperref
\chapter*{Agradecimientos} % Opción con * para que no aparezca en TOC ni numerada
\addcontentsline{toc}{chapter}{Agradecimientos} % Añade al TOC.

 Aunque es un apartado opcional, haremos bueno el refrán \emph{<<es de bien nacidos, ser agradecidos>>} si empleamos este espacio como un medio para agradecer a todos los que, de un modo u otro, han hecho posible que el trabajo realizado \emph{llegue a buen puerto}. Esta sección es ideal para agradecer a directores, profesores, mentores, familiares, compañeros, amigos, etc. 
 
 Estos agradecimientos pueden ser tan personales como se desee e incluir anécdotas y chascarrillos, pero \emph{nunca deberían ocupar más de una página}.

\makeatletter		
\begin{flushright}
	\vspace{1,5cm}
	\textit{\@autor}\\
	\@cityTF, \@yearTF
\end{flushright}
\makeatother % Agradecimientos etc.
% -------------------------
%
% ÍNDICES: 
% EDITAR: Si alguno de los índices no existe, su inclusión se puede comentar.
%
% -------------------------
\setindexnames % Ajusta nombres (sólo en español).
\pagestyle{fancy} % Estilo de página ajustado por fancyhdr

%--- Índice general
\cleardoublepage
\phantomsection % Necesario con hyperref
\pdfbookmark[0]{Índice general}{idx_toc}% idx_toc.0 % Bookmark en PDF
\tableofcontents  % Índice general
% Todos los listados se han incluido en el índice gral. de contenidos. De 
%modo automático también quedan añadidos a los bookmarks del PDF. Si se 
%desean 
%eliminiar del TOC se pueden comentar el comando \addcontensline.
%---

%--- Índice de figuras
\cleardoublepage
\phantomsection % Necesario con hyperref
\addcontentsline{toc}{chapter}{\listfigurename} % Añade la lista de figuras al TOC (también a bookmarks en PDF)
%\pdfbookmark[0]{\listfigurename}{idx_lof}% idx_lof.0 % Bookmark en PDF
\listoffigures    % Índice de figuras (opcional)
%---

%--- Índice de tablas
\cleardoublepage
\phantomsection % Necesario con hyperref
\addcontentsline{toc}{chapter}{\listtablename} % Añade la lista de tablas al TOC (también a bookmarks en PDF)
%\pdfbookmark[0]{\listtablename}{idx_lot}% idx_lot.0 % Bookmark en PDF
\listoftables % Índice de tablas (opcional)
%---
 % Índice de contenido, figuras, tablas
%--- (FIN FRONTMATTER)



% -------------------------
% -------------------------
% -------------------------
% -------------------------
%
%--- MAINMATTER
% Capítulos del documento
\savepagecnt % Se almacena en un contador interno el nº de páginas actual
% Debe ir antes de \mainmatter (antes de que se reinicie el cnt page)
\mainmatter
% Justo antes del primer capítulo del libro. Activa la numeración con números arábigos y reinicia el contador de páginas.

% Ajusta valor de cabeceras y pies a comienzo de capítulo
\ifpageonfooter
\else
	\cleanhdfirst
\fi

% Reajuste del n. de pág. consecutivo para no reiniciar pag. en Cáp. 1
%\contpagination % Comentado para reiniciar paginación (pag. 1)
% -------------------------
% --- CAPÍTULOS: Un fichero por capítulo.
% -------------------------
% OPT: Si deseas modificar el interlineado de una línea (por defecto)
% puedes emplear los comandos señalados (proporcionados por paquete `setspace')
%\onehalfspacing % Ajusta interlineado a 1.5 líneas


\chapter{Introducción}
\label{cap:Introduccion}

En los estudios de ingeniería y publicaciones relacionadas, el estilo de citación más empleado es el numérico. En este caso se cita una referencia mediante un número incluido entre corchetes (p.~ej., [\#]) que indica el orden de la referencia citada en la lista de referencias incluida en la sección correspondiente. Con este tipo de citación se suele utilizar dos tipos de ordenación para la lista de referencias. Estos son: el orden alfabético de las referencias\footnote{Este es el estilo por defecto o \texttt{plain} empleado con \LaTeX.} y el orden de citación de la referencia en el texto\footnote{Obtenido con los estilos \texttt{unsrt} y \texttt{ieeetr}, entre otros.}. Por el contrario, en las ciencias naturales y sociales es más habitual un estilo de citación autor-año como el recogido en las normas APA (\href{https://apastyle.apa.org/learn/faqs/format-bibliography}{American Pshychological Association}).

Con \LaTeX{} existen varias alternativas para utilizar el estilo de citación autor-año. En este documento se propone una de las más sencillas. Esta consiste en el empleo del paquete \href{https://ctan.javinator9889.com/biblio/bibtex/contrib/apacite/apacite.pdf}{\texttt{apacite}} con la opción \texttt{natbibapa} que proporciona un estilo de citación siguiendo el formato sugerido por APA y los comandos de citación del paquete \texttt{natbib} (cargado de modo implícito). Con estas opciones la primera citación de una referencia incluye a todos los autores \citep[como por ejemplo en][]{oetiker06} y el resto solo una lista corta \citep{oetiker06}. Cuando existe una citación múltiple, las citas aparecen en el mismo orden que en la lista de referencias \citep[como por ejemplo en][]{lamport94, cascales00,goos04,kopka04}.

\section{Citación con \texttt{apacite} y \texttt{natbib}}
Con el paquete \texttt{apacite} y la opción \texttt{natbibapa} se proporcionan varios comandos de citación que admiten dos opciones para añadir texto previo (prefijo) a la citación y posterior (sufijo) a la misma, como en:

   \verb+\citep[ver][página 11]{salido15}+: \citep[ver][página 5]{salido15} 

\newpage
Los comandos de citación más frecuentes se resumen en la lista siguiente:

\begin{itemize}
    \item \verb+\citep+: Citación entre paréntesis.
    \item \verb+\citet+: Citación textual (sin paréntesis).
    \item \verb+\citep*+: Citación con lista completa de autores entre paréntesis.
    \item \verb+\citet*+: Citación textual con lista completa de autores (sin paréntesis).
    \item \verb+\citeauthor+: Cita solo el autor (sin paréntesis).
    \item \verb+\citeyear+: Cita solo el año (sin paréntesis).
    \item \verb+\citeyearpar+: Cita solo el año (con paréntesis).
\end{itemize}

\chapter{Objetivo}
\label{cap:Objetivo}

Introduce y motiva la problemática (i.e.\emph{\ ¿\textsc{cuál} es el problema que se plantea y \textsc{por qué} es interesante su resolución?})


\chapter{Resultados}
\label{cap:Resultados}


\chapter{Conclusiones}
\label{cap:Conclusiones}

En este capítulo se realizará un juicio crítico y discusión sobre los resultados obtenidos. \emph{Cuidado, esta discusión no debe confundirse con una valoración del enriquecimiento personal que supone la realización del trabajo como culminación de una etapa académica}. Aunque de gran importancia, esta última valoración debe quedar fuera de la memoria del trabajo y solo debe ahondarse en ella ante requerimiento explícito del comité en el acto de defensa.

Si es pertinente deberá incluir información sobre trabajos derivados como publicaciones o ponencias en preparación, así como trabajos futuros \emph{(solo si estos están iniciados o planificados en el momento que se redacta el texto)}. Evitar hacer una lista de posibles mejoras. Contrariamente a lo que alguno pueda pensar generalmente aportan impresión de trabajo incompleto o inacabado.\footnote{Puede reflexionarse en ello por si en la defensa del trabajo se pregunta sobre estas posibles mejoras.}









%--- (FIN FRONTMATTER)







% -------------------------
% No olvides retornar al interlineado sencillo en el resto del documento.
\singlespacing
% -------------------------
% -------------------------
% -------------------------
% -------------------------
%
%--- BACKMATTER
%BEGIN_FOLD
% \backmatter 
% Comentado para que apéndices aparezcan numerados después de la bibliografía


% -------------------------
% --- BIBLIOGRAFÍA
% -------------------------
\cleardoublepage % Necesario para ajustar el avance de página
\phantomsection  % Ojo necesario con hyperref.
\addcontentsline{toc}{chapter}{\bibname} % Añade la bibliografía al Índice de contenidos

% EDITAR: Definición del estilo empleado en la bibliografía
% Comentar si se emplea lista de \bibitem{citekey}
\bibliographystyle{apacite}
% apacite:   No numérica, con referencias ordenadas alfabéticamente por apellido %            de autor.
%---
% Comentar si se emplea lista de \bibitem{citekey} 
\bibliography{biblio} % Nombre del fichero .bib (sin extensión)
% -------------------------

% -------------------------
% Sustituir por entorno 
% Contenido obtenido de fichero auxiliar bbl donde de puede editar
% manualmente cada una de las entradas.
%\begin{thebibliography}{10}
%\bibitem{borbon21}
%Alexánder Borbón and Walter Mora.
%\newblock {\em Edición de textos científicos con {\LaTeX{}}. {C}omposición,
%  diseño editorial, gráficos, {I}nkscape, {T}ikz y presentaciones {B}eamer}.
%\newblock Instituto Tecnológico de Costa Rica, 2 edition, 2021.
%
%\bibitem{cascales03}
%Bernardo Cascales, Pascual Lucas, José~M. Mira, Antonio~J. Pallarés, and
%  Salvador {Sánchez-Pedreño}.
%\newblock {\em El libro de {\LaTeX}}.
%\newblock Pearson Education, 2005.
%\newblock {ISBN}: 9788420537795.
%
%\bibitem{esi19}
%{Escuela Superior de Informática, Universidad de Castilla-La Mancha}.
%\newblock Guía de estilo y formato para {T}rabajos {F}in de {G}rado.
%\newblock {URL:}
%  \url{https://pruebasaluuclm.sharepoint.com/sites/esicr/tfg/SiteAssets/SitePages/Inicio/20190304-GuiaEstiloFormatoTFG.pdf},
%  March 2019.
%\newblock Último acceso: sep. 2021.
%
%\bibitem{goos07}
%Michel Goossens, Frank Mittelbach, Sebastian Rahtz, Denis Roeget, and Herbert
%  Vo{\ss}.
%\newblock {\em The {\LaTeX{}} graphics companion}.
%\newblock Addison-Wesley Reading, MA, second edition, 2007.
%\newblock {ISBN:} 9780321508928.
%
%\bibitem{grat99}
%George~A. Grätzer.
%\newblock {\em First steps in {\LaTeX}}.
%\newblock Springer Verlag, October 1999.
%\newblock {ISBN:} 0817641327.
%
%\bibitem{grat07}
%George~A. Grätzer.
%\newblock {\em More math into {\LaTeX}}.
%\newblock Birkhauser, fourth edition, 2007.
%\newblock {ISBN}: 978-3319237954.
%
%\bibitem{lamport94}
%Leslie Lamport.
%\newblock {\em {\LaTeX}: A document preparation system}.
%\newblock Addison-Wesley, second edition, June 1994.
%\newblock {ISBN}: 978-0201529838.
%
%\bibitem{mittelbach04}
%Frank Mittelbach and Michel Goossens.
%\newblock {\em The {\LaTeX{}} companion}.
%\newblock Addison, Reading, Mass, second edition, 2004.
%\newblock {ISBN:} 0-201-36299-6.
%
%\bibitem{oetiker14}
%Tobias Oetiker, Hubert Partl, Irene Hyna, and Elisabeth Schlegl.
%\newblock {\em La introducción no-tan-corta a {\LaTeX2e}}, 2014.
%
%\bibitem{salido10}
%Jesús Salido.
%\newblock Curso: {\LaTeX{}} esencial para preparación de tfg, tesis y otros
%  documentos académicos.
%\newblock {URL:} \url{http://visilab.etsii.uclm.es/?page_id=1468}, 2010.
%\newblock Último acceso: sep. 2021.
%
%\bibitem{salido19}
%Jesús Salido.
%\newblock {Plantilla LaTeX para Trabajo Fin de Estudios en Ingeniería
%  Informática - UCLM (España)}.
%\newblock Zenodo, DOI: \url{https://doi.org/10.5281/zenodo.4561708}, 2019.
%
%\bibitem{uclm20}
%{Servicio de publicaciones UCLM}.
%\newblock Propiedad intelectual.
%\newblock {D}ocumentos elaborados por el {G}rupo de {G}estión del
%  {C}onocimiento y {P}ropiedad {I}ntelectual de la {Universidad de Castilla-La
%  Mancha}. {URL:} \url{https://publicaciones.uclm.es/propiedad-intelectual/}.
%\newblock Último acceso: jun. 2022.
%
%\bibitem{uc3m21}
%{Universidad Carlos III de Madrid}.
%\newblock Guía temática del {TFG}.
%\newblock {URL:} \url{https://uc3m.libguides.com/TFG}, 2021.
%\newblock Último acceso: sep-2021.
%
%\bibitem{unican18}
%{Universidad de Cantabria}.
%\newblock Cómo usar imágenes en trabajos.
%\newblock Artículo técnico disponible en {URL:}
%  \url{https://web.unican.es/buc/Documents/Formacion/guia_imagenes.pdf}, 2018.
%\newblock Último acceso: sep. 2021.
%
%\bibitem{uoc}
%{Universitat Oberta de Catalunya}.
%\newblock Comunicación eficaz y redacción.
%\newblock {URL:}
%  \url{https://www.uoc.edu/portal/es/servei-linguistic/redaccio/10-recomanacions/index.html}.
%\newblock Último acceso: sep. 2021.
%
%\bibitem{Vijayasarathy_2016}
%Leo~R. Vijayasarathy and Charles~W. Butler.
%\newblock {Choice of Software Development Methodologies: Do Organizational,
%  Project, and Team Characteristics Matter?}
%\newblock {\em {IEEE} Software}, 33(5):86--94, sep 2016.
%
%\bibitem{wikibookLaTex10}
%WikiMedia.
%\newblock {\LaTeX{} Wikibook}.
%\newblock {URL:} \url{http://en.wikibooks.org/wiki/LaTeX}, 2010.
%\newblock Último acceso: sep. 2021.
%
%\bibitem{zozaya17}
%Leonor Zozaya.
%\newblock Redacción de textos. recomendaciones para presentar trabajos
%  académicos.
%\newblock Blog disponible en {URL:} \url{http://redaccion.hypotheses.org/},
%  2017.
%\newblock {ISSN}: 2444-8885.
%\end{thebibliography}
  





% -------------------------
% - ANEXOS: Comentar si no se desean incluir. [OPT.]
% - Mover si se desea que aparezcan antes de la bibliografía.
% -------------------------
\appendix
\ifspanish
	\part*{\sffamily ANEXOS}
\else
	\part*{\sffamily APPENDICES}
\fi
% Tras este punto los capítulos se numeran con letras.
% Aquí todos los apéndices necesarios
\chapter{El primer anexo}
\label{cap:AnexoA}

Los anexos se incluirá de modo opcional material suplementario que podrá consistir en breves manuales, listados de código fuente, esquemas, planos y en general aquello que complementa a la memoria, pero no se incluye en el texto principal. Se recomienda que no sean excesivamente voluminosos, aunque su extensión no está sometida a la regulación por normativa, ya que esta afecta únicamente al texto principal de la memoria. 




 % Apéndice A (opcionales)

%END_FOLD
%--- (FIN DOCUMENTO)
\end{document}